% Л1. Платформа .NET и dotnet CLI -- слайды.
% Спутник текста лекции в ../L01-platform.md. Текст -- источник истины:
% слайды несут код, схемы и короткие утверждения, а подача живёт в \note{}.
%
% Сборка:  make          (latexmk -pdf; движок -- pdflatex)
%          make notes    (то же плюс заметки докладчика на втором экране)
%
% Дом. правила -- в ../CLAUDE.md, тема -- в ../../../../theme/.

\documentclass[aspectratio=169,10pt,t]{beamer}

\usetheme{dotnet}
\usetikzlibrary{arrows.meta,positioning,fit}

\dotnetlecture{Л1}

% Заметки докладчика скрыты по умолчанию; `make notes` собирает версию
% для второго экрана.
\ifdefined\shownotes
  \setbeameroption{show notes on second screen=right}
\fi

\title[Л1. Платформа .NET]{Платформа .NET и dotnet CLI}
\author{Дмитрий Курочкин}
\date{}

\begin{document}

% ------------------------------------------------------------------ титул --
\begin{frame}[plain]
  \titlepagednc{Лекция 1}{Платформа .NET и dotnet CLI}{%
    Что происходит между вашим текстом и процессом в операционной системе}{%
    Дмитрий Курочкин\\[0.35em]
    {\color{muted}\small\texttt{github.com/kudima03}}}
\end{frame}
\note{80 минут. Начинаем со знакомства, потом сразу вопрос в зал.}

% ==================================================== 0. ЗНАКОМСТВО ========
\section{0. Знакомство и структура}

\begin{frame}[plain]
  \sectionopener{0}{Знакомство, дисклеймеры и структура}{Кто я, как устроен курс
    и о чём нельзя не сказать в 2026 году.}
\end{frame}
\note{Бюджет: 9 минут.}

\begin{frame}{О чём договоримся на входе}
  \slidetop
  \claim{Три вещи, которые определяют всё остальное в этом курсе.}

  \begin{itemize}\itemsep0.7em
    \item Я не преподаватель, а \hi{практик} --- приглашённый спикер,
          который пишет на .NET каждый день.
    \item ИИ требует мастера: инструмент не снимает с вас последний контроль.
    \item На выходе --- рабочий сервис, собранный своими руками.
  \end{itemize}
\end{frame}
\note{Не зачитывать список. Первый пункт -- переход к следующему слайду.}

\begin{frame}{Коротко о себе}
  \slidetop
  \claim{.NET --- моя основная технология, под которую я пишу каждый день.}

  \begin{minipage}[t]{0.60\textwidth}
    \eyebrow{чем занимаюсь}\par\vspace{0.6em}
    Около года удерживаю первое место по версии
    \hi{\texttt{commiters.top}}.\\[0.5em]
    Автор и главный maintainer экосистемы \texttt{Pure} --- по ней есть
    отдельный доклад.
  \end{minipage}\hfill
  \begin{minipage}[t]{0.32\textwidth}
    \eyebrow{контакты}\par\vspace{0.6em}
    % Заменить рамку на \includegraphics{qr.pdf}, когда QR будет готов.
    \begin{tikzpicture}
      \draw[hairline,line width=0.6pt] (0,0) rectangle (2.2,2.2);
      \node[muted,font=\scriptsize] at (1.1,1.1) {QR};
    \end{tikzpicture}\\[0.4em]
    {\color{muted}\scriptsize GitHub \textperiodcentered{} почта
      \textperiodcentered{} телеграм}
  \end{minipage}
\end{frame}
\note{Ссылки -- на слайде, вслух только суть. Дальше сразу вопрос про
  дотнетчиков.}

\begin{frame}[plain]
  \statement{\hi{КРИТИКУЙ!}}
\end{frame}
\note{Курс в первую очередь для дотнетчиков в зале. Обратная связь нужна
  прежде всего от них, прямо посреди лекции. Это единственное слово во весь
  экран за всю лекцию -- оно и должно запомниться.}

\begin{frame}{.NET проектировали люди, которые хорошо знали Java}
  \slidetop
  \claim{При создании .NET много было взято из Java как из \hi{референса} ---
    и хорошего, и плохого.}

  \begin{center}
    \begin{tikzpicture}[x=3.05cm,font=\small]
      \draw[hairline,line width=0.6pt] (-0.15,0) -- (3.35,0);
      \foreach \i/\name/\year in {0/Turbo Pascal/1983, 1/Delphi/1995,
                                  2/Visual J++/1996, 3/C\#/2000} {
        \fill[muted] (\i,0) circle (1.6pt);
        \node[above=0.45em,ink] at (\i,0) {\name};
        \node[below=0.5em,muted,font=\scriptsize] at (\i,0) {\year};
      }
      \node[accent,font=\scriptsize,above=2.1em] at (2,0) {Java, изнутри};
    \end{tikzpicture}
  \end{center}

  \slidegap
  \aside{Андерс Хейлсберг вёл Visual J++ до того, как спроектировал C\#.}
\end{frame}
\note{Фраза товарища про запах -- голосом, не на слайде. Вывод: язык вышел
  свежее, менее многословным, с более чистыми подходами.}

\begin{frame}{Структура курса}
  \slidetop
  \claim{15 лекций по 80 минут и 15 лабораторных, \hi{след в след}.}

  \eyebrow{примерно с лабы 8 растёт один сервис}\par
  \vspace{0.8em}
  \begin{center}
    \begin{tikzpicture}[font=\small,node distance=0.32cm]
      \node (a) [muted] {контракт};
      \node (b) [right=of a] {$\rightarrow$};
      \node (c) [right=of b] {настройки};
      \node (d) [right=of c] {$\rightarrow$};
      \node (e) [right=of d] {безопасность};
      \node (f) [right=of e] {$\rightarrow$};
      \node (g) [right=of f] {кэш и логи};
      \node (h) [right=of g] {$\rightarrow$};
      \node (i) [right=of h] {база};
      \node (j) [right=of i] {$\rightarrow$};
      \node (k) [right=of j,accent] {тесты};
    \end{tikzpicture}
  \end{center}

  \slidegap
  \aside{Первые лабы маленькие и отдельные. Дальше вы растите один сервис
    слой за слоем и понимаете в нём каждую строчку.}
\end{frame}
\note{Обещание с моей стороны: не рассказывать то, что вы и так знаете.
  Просьба с вашей: спрашивайте, особенно неправильное.}

% ================================================= 1. ЧТО ТАКОЕ .NET =======
\section{1. Что такое .NET}

\begin{frame}[plain]
  \sectionopener{1}{Что такое .NET}{Четыре части платформы, четыре аббревиатуры,
    версии и разница между SDK и Runtime.}
\end{frame}
\note{Бюджет: 13 минут.}

\begin{frame}[plain]
  \vote{.NET --- это язык программирования?}
\end{frame}
\note{Разогрев. Большинство скажет <<нет>> и будет право. Дальше -- почему
  вопрос вообще стоит задавать.}

\begin{frame}{Платформа, а не язык}
  \slidetop
  \claim{В .NET четыре части. Слова \hi{<<C\#>>} среди них нет.}

  \begin{center}
    \begin{tikzpicture}[font=\small]
      \node[muted,font=\scriptsize] at (0,2.45)
        {C\#\quad F\#\quad VB\quad C++ via CLR\quad\ldots{} порядка 70 языков};
      \draw[hairline,line width=0.6pt] (-4.6,2.15) -- (4.6,2.15);
      \foreach \x/\t in {-3.45/рантайм, -1.15/{стандартная\\библиотека},
                         1.15/инструменты, 3.45/{экосистема\\пакетов}} {
        \node[draw=hairline,line width=0.6pt,fill=codebg,
              minimum width=2.05cm,minimum height=1.25cm,align=center,
              font=\small] at (\x,1.2) {\t};
      }
      \node[accent,font=\scriptsize] at (0,0.1)
        {свой язык? напиши компилятор};
    \end{tikzpicture}
  \end{center}
\end{frame}
\note{Рантайму всё равно, на чём вы написали. Ему важно, во что это
  скомпилировалось.}

\begin{frame}{Четыре аббревиатуры}
  \slidetop
  \claim{Расшифровываю один раз. Дальше пользуюсь без расшифровки.}

  \begin{tabular}{@{}p{1.7cm}p{9.5cm}@{}}
    \textbf{CLR} & Common Language Runtime --- \hi{исполняет}: загружает код,
                   компилирует, управляет памятью \\[0.55em]
    \textbf{CTS} & Common Type System --- определяет, что такое тип; поэтому
                   класс из F\# виден из C\# \\[0.55em]
    \textbf{CLS} & Common Language Specification --- договор о совместимости
                   между языками \\[0.55em]
    \textbf{BCL} & Base Class Library --- базовый минимальный набор
                   компонентов, включая примитивы \\
  \end{tabular}

  \slidegap
  \aside{CTS --- это стандарт вилки: чайник одного производителя и лампа
    другого работают от одной сети.}
\end{frame}
\note{Момент запомнить. CLS на практике почти не всплывает -- одной фразой.}

\begin{frame}{Версии: раз в год, каждый ноябрь}
  \slidetop
  \claim{Чётные номера --- \hi{LTS}, поддержка три года. Нечётные --- STS.}

  \begin{center}
    \begin{tikzpicture}[x=3.05cm,font=\small]
      \draw[hairline,line width=0.6pt] (-0.2,0) -- (3.4,0);
      \foreach \i/\v/\k/\y in {0/.NET 8/LTS/2023, 1/.NET 9/STS/2024,
                               2/.NET 10/LTS/2025, 3/.NET 11/STS/2026} {
        \fill[muted] (\i,0) circle (1.6pt);
        \node[above=0.45em] at (\i,0) {\v};
        \node[below=0.5em,muted,font=\scriptsize] at (\i,0) {\k \textperiodcentered{} \y};
      }
      \node[accent,font=\small,above=1.55em] at (2,0) {мы здесь};
    \end{tikzpicture}
  \end{center}

  \slidegap
  \aside{Новый проект --- на последней LTS, если нет причины, которую можно
    произнести вслух на совещании. Bump до новой версии --- через месяц-два
    после выхода.}
\end{frame}
\note{Держите проект так, чтобы переезд был сменой одной версии, без
  переписывания кода.}

\begin{frame}{SDK и Runtime: что ставить}
  \slidetop
  \claim{\hi{SDK} себе, \hi{Runtime} туда, где живёт приложение.}

  \versus{SDK}{%
    Компилятор, CLI, шаблоны проектов, MSBuild.\\[0.5em]
    И рантайм внутри --- чтобы запускать у себя.\\[0.5em]
    {\color{muted}Ставит разработчик.}}{Runtime}{%
    Только то, что нужно, чтобы запустить готовое приложение.\\[0.5em]
    Компилятора нет. Шаблонов нет.\\[0.5em]
    {\color{muted}Стоит на сервере или у пользователя.}}

  \slidegap
  \aside{Номера версий у них похожие, но не одинаковые: SDK обновляется чаще
    и сам по себе.}
\end{frame}
\note{Что стоит на сервере -- вы обязаны понимать. Модели поставки -- большой
  кусок следующей лекции.}

\begin{frame}[fragile]{Что стоит на этой машине}
  \slidetop
\begin{lstlisting}[style=term]
$ dotnet --info

.NET SDKs installed:
  8.0.404  [/usr/share/dotnet/sdk]
  10.0.100 [/usr/share/dotnet/sdk]

.NET runtimes installed:
  Microsoft.AspNetCore.App 8.0.11  [/usr/share/dotnet/shared]
  Microsoft.AspNetCore.App 10.0.0  [/usr/share/dotnet/shared]
  Microsoft.NETCore.App    8.0.11  [/usr/share/dotnet/shared]
  Microsoft.NETCore.App    10.0.0  [/usr/share/dotnet/shared]
\end{lstlisting}
  \vspace{0.7em}
  \aside{Рантаймов несколько, и они двух видов: базовый и тот, что для веба.
    Каждое приложение при старте выбирает свой --- по файлу, который мы
    увидим в конце лекции.}
\end{frame}
\note{Если у кого-то команда не выводит ничего похожего -- SDK нет, до лабы
  поставить.}

\begin{frame}[plain]
  \keyline{.NET --- это платформа.\\ C\# --- только вход в неё.}
\end{frame}
\note{Пауза на вопросы по разделу.}

% ================================================== 2. ПУТЬ КОДА ==========
\section{2. Путь кода}

\begin{frame}[plain]
  \sectionopener{2}{Путь кода}{Исходник $\rightarrow$ IL $\rightarrow$ машинный
    код: два компилятора, два этапа и всё, что из этого следует.}
\end{frame}
\note{Бюджет: 20 минут. Главный раздел лекции.}

\begin{frame}[plain]
  \vote{.NET компилируется или интерпретируется?}
\end{frame}
\note{Голосование в три руки. Правы все трое -- но пойдём по порядку.}

\begin{frame}[fragile]{Что на самом деле лежит в dll}
  \slidetop
  \begin{minipage}[t]{0.46\textwidth}
    \eyebrow{исходник}\par\vspace{0.5em}
\begin{lstlisting}[style=csharp]
static int Add(int a, int b)
    => a + b;
\end{lstlisting}
  \end{minipage}\hfill
  \begin{minipage}[t]{0.50\textwidth}
    \eyebrow{то же самое в IL}\par\vspace{0.5em}
\begin{lstlisting}[style=plain]
.method static int32
        Add(int32 a, int32 b) cil managed
{
    .maxstack 8
    IL_0000: ldarg.0
    IL_0001: ldarg.1
    IL_0002: add
    IL_0003: ret
}
\end{lstlisting}
  \end{minipage}

  \vspace{0.9em}
  \aside{Ни слова о процессоре \textperiodcentered{} сохранены имена и типы
    --- это метаданные \textperiodcentered{} читается глазами, поэтому
    \hi{секретов там быть не должно}.}
\end{frame}
\note{Не учить IL. Увидеть три вещи. На метаданных держится половина
  платформы: отладчик, рефлексия, сериализация.}

\begin{frame}{Компилируется дважды}
  \slidetop
  \claim{Спор <<компилируется или интерпретируется>> бессмысленный.
    Правильный ответ --- \hi{дважды}.}

  \begin{center}
    \begin{tikzpicture}[font=\small,
      box/.style={draw=hairline,line width=0.6pt,fill=codebg,
                  minimum height=0.95cm,align=center,inner sep=0.45em},
      ar/.style={-{Straight Barb[length=1.6mm]},muted}]
      \node[box] (src) {Program.cs};
      \node[box,right=1.15cm of src] (il) {dll с IL};
      \node[box,right=1.15cm of il] (mc) {машинный код\\в памяти};
      \draw[ar] (src) -- node[above,font=\scriptsize,muted]{Roslyn} (il);
      \draw[ar] (il) -- node[above,font=\scriptsize,muted]{JIT} (mc);
      \node[muted,font=\scriptsize,below=0.55em] at (src.south)
        {ваша машина, один раз};
      \node[muted,font=\scriptsize,below=0.55em,align=center] at (mc.south)
        {машина пользователя,\\при первом вызове метода};
    \end{tikzpicture}
  \end{center}

  \slidegap
  \aside{Артефакт на диске один: dll с IL. Машинный код рождается заново на
    каждом запуске и умирает вместе с процессом.}
\end{frame}
\note{Те, кто голосовал за <<интерпретируется>>, почувствовали правильно, что
  что-то происходит в рантайме.}

\begin{frame}{Чем мы за это платим}
  \slidetop
  \claim{Два больших плюса и один минус.}

  \begin{itemize}\itemsep0.75em
    \item[\textcolor{accent}{$+$}] Артефакт не зависит от платформы: одна dll
          работает везде, где есть рантайм.
    \item[\textcolor{accent}{$+$}] JIT знает конкретную машину и её набор
          инструкций --- то, чего не знает компилятор C++ на вашем ноутбуке.
    \item[\textcolor{warn}{$-$}] \hi{Прогрев на старте}: каждый метод при
          первом вызове надо скомпилировать.
  \end{itemize}
\end{frame}
\note{Здесь я упрощаю -- дотнетчики, поправьте: JIT не <<подстраивается под
  всё>>, речь про набор инструкций. И у C++ есть -march=native.}

\begin{frame}{Многоуровневая компиляция}
  \slidetop
  \claim{Рантайм не выбирает между <<быстро и плохо>> и <<медленно и
    хорошо>>. Он делает \hi{и то и другое}, по очереди.}

  \begin{center}
    \begin{tikzpicture}[font=\small,
      box/.style={draw=hairline,line width=0.6pt,fill=codebg,
                  minimum height=0.95cm,align=center,inner sep=0.45em},
      ar/.style={-{Straight Barb[length=1.6mm]},muted}]
      \node[box] (t0) {уровень 0\\быстро и грубо};
      \node[box,right=1.9cm of t0] (t1) {уровень 1\\полная оптимизация};
      \draw[ar] (t0) -- node[above,font=\scriptsize,muted]{счётчик вызовов}
                       node[below,font=\scriptsize,muted]{в фоне, с подменой кода} (t1);
      \node[muted,font=\scriptsize,below=0.55em] at (t0.south)
        {при первом вызове};
      \node[muted,font=\scriptsize,below=0.55em] at (t1.south)
        {когда метод стал горячим};
    \end{tikzpicture}
  \end{center}

  \slidegap
  \aside{Первый запрос после старта медленнее второго --- это не баг.
    Прогрев --- свойство процесса: перезапустили, и всё начинается заново.}
\end{frame}
\note{Длинный цикл тоже горячий, хотя счётчик вызовов равен единице.}

\begin{frame}{Debug и Release: что реально меняется}
  \slidetop
  \begin{tabular}{@{}p{5.0cm}p{2.9cm}p{2.9cm}@{}}
     & \eyebrow{debug} & \eyebrow{release} \\[0.5em]
    Оптимизации компилятора   & нет        & да \\[0.45em]
    Оптимизации JIT           & запрещены  & разрешены \\[0.45em]
    Константа \texttt{DEBUG}  & есть       & нет \\[0.45em]
    Файл \texttt{pdb}         & есть       & \hi{есть} \\[0.45em]
    Назначение                & отладка    & измерение и поставка \\
  \end{tabular}

  \slidegap
  \aside{Release означает <<с оптимизациями>>, а не <<без символов>>. Символы
    нужны ради номеров строк в трассировке исключения.}
\end{frame}
\note{В Debug в сборку кладётся прямая пометка для JIT: не оптимизируй тоже.}

\begin{frame}[plain]
  \statement{Бенчмарк, снятый в \hi{Debug}, не значит ничего.}
\end{frame}
\note{Сборка по умолчанию -- Debug. Когда вам покажут число, первый вопрос:
  в какой конфигурации?}

\begin{frame}[plain]
  \keyline{Ваш код компилируется дважды:\\ в IL на вашей машине, в машинный
    код --- на той, где он запускается.}
\end{frame}
\note{Всё остальное, от прогрева до Release, следует из этого.}

% ============================================ 3. ПРОЕКТ И РЕШЕНИЕ =========
\section{3. Проект и решение}

\begin{frame}[plain]
  \sectionopener{3}{Проект и решение}{csproj --- это манифест для MSBuild,
    а не файл настроек.}
\end{frame}
\note{Бюджет: 15 минут.}

\begin{frame}[fragile]{Файл проекта целиком}
  \slidetop
\begin{lstlisting}[style=xml]
<Project Sdk="Microsoft.NET.Sdk">

  <PropertyGroup>
    <OutputType>Exe</OutputType>
    <TargetFramework>net10.0</TargetFramework>
    <ImplicitUsings>enable</ImplicitUsings>
    <Nullable>enable</Nullable>
  </PropertyGroup>

</Project>
\end{lstlisting}
  \vspace{0.7em}
  \aside{Девять строк, и это всё. Файл короткий потому, что правила лежат
    в \hi{SDK}, а здесь только отклонения от умолчаний.}
\end{frame}
\note{Разбирать строку за строкой. Лет десять назад каждый файл надо было
  перечислять руками.}

\begin{frame}[fragile]{Два способа подключить чужой код}
  \slidetop
\begin{lstlisting}[style=xml]
<ItemGroup>
  <PackageReference Include="Newtonsoft.Json" Version="13.0.3" />
  <ProjectReference Include="../Shared/Shared.csproj" />
</ItemGroup>
\end{lstlisting}
  \vspace{0.9em}
  \versus{PackageReference}{%
    Из реестра NuGet, \hi{замороженная} версия.\\[0.4em]
    {\color{muted}Чтобы получить изменения, нужен новый пакет.}}%
  {ProjectReference}{%
    Соседний проект в том же репозитории.\\[0.4em]
    {\color{muted}Собирается вместе и подхватывает изменения.}}

  \slidegap
  \aside{Свой код из того же репозитория --- проектом. Чужой и выпускаемый
    по версиям --- пакетом.}
\end{frame}
\note{Версия должна быть зафиксирована -- здесь или централизованно. Здесь
  я упрощаю: без Version тоже соберётся.}

\begin{frame}{Решение --- это список, а не контейнер}
  \slidetop
  \claim{Проект живёт сам по себе. Решение просто \hi{перечисляет} проекты.}

  \begin{center}
    \begin{tikzpicture}[font=\small,
      box/.style={draw=hairline,line width=0.6pt,fill=codebg,
                  minimum height=0.8cm,minimum width=2.5cm,align=center,
                  inner sep=0.4em},
      ar/.style={-{Straight Barb[length=1.6mm]},muted}]
      \node[box] (sln) {\texttt{Product.sln}};
      \node[box,right=2.2cm of sln,yshift=1.05cm]  (p1) {\texttt{Api.csproj}};
      \node[box,right=2.2cm of sln]                (p2) {\texttt{Domain.csproj}};
      \node[box,right=2.2cm of sln,yshift=-1.05cm] (p3) {\texttt{Tests.csproj}};
      \draw[ar] (sln.east) -- (p1.west);
      \draw[ar] (sln.east) -- (p2.west);
      \draw[ar] (sln.east) -- (p3.west);
    \end{tikzpicture}
  \end{center}

  \slidegap
  \aside{Что это даёт: группировку и правильный порядок сборки. Нужно, как
    только проектов стало больше одного.}
\end{frame}
\note{IDE старательно поддерживает недоразумение, что проекты живут внутри
  решения. Сегодня соберём проект вообще без решения.}

\begin{frame}{Что появляется после сборки}
  \slidetop
  \versus{obj --- черновики}{%
    \texttt{project.assets.json}, промежуточные файлы, сгенерированные
    \texttt{using}.\\[0.5em]
    {\color{muted}Кухня.}}{bin/Debug/net10.0 --- результат}{%
    dll и всё, что нужно рядом, чтобы приложение запустилось.\\[0.5em]
    {\color{muted}Окно выдачи.}}

  \slidegap
  \hrulethin{0.4}
  \rulegap

  {\color{warn}\large bin и obj не коммитятся в репозиторий. Никогда.}

  \slidegap
  \aside{Это производное от исходников: всё получается заново одной командой.
    Если сборка ведёт себя странно --- удаляйте и пересобирайте.}
\end{frame}
\note{В obj лежат пути с вашей машины, и у коллеги они не совпадут.}

\begin{frame}[plain]
  \keyline{csproj --- манифест сборки, который вы обязаны уметь читать.\\
    Всё, что порождает сборка, живёт в bin и obj.}
\end{frame}
\note{И не имеет ценности само по себе.}

% ================================================== 4. DOTNET CLI =========
\section{4. dotnet CLI}

\begin{frame}[plain]
  \sectionopener{4}{dotnet CLI}{Единственный интерфейс к платформе, который
    существует везде, где живёт ваше приложение.}
\end{frame}
\note{Бюджет: 18 минут.}

\begin{frame}{Почему терминал, а не кнопка}
  \slidetop
  \claim{Объяснение не про <<так правильно>>, а про конкретную выгоду.}

  \begin{enumerate}\itemsep0.7em
    \item На сервере сборки, в CI и в контейнере \hi{IDE нет}. Там есть
          ровно эти команды.
    \item IDE --- надстройка над ними же. Терминал показывает ошибку целиком,
          а не её сокращённую версию в окошке.
    \item Команду можно записать в файл, положить в скрипт и передать
          коллеге. Последовательность кликов --- нельзя.
  \end{enumerate}

  \slidegap
  \aside{За кнопкой <<запустить>> прячется пять действий. Вы видите одно.}
\end{frame}
\note{Аналогия с механикой и автоматом. LLM, кстати, запускает ровно эти же
  команды -- ничего другого у неё нет.}

\begin{frame}[fragile]{Всё начинается с шаблона}
  \slidetop
\begin{lstlisting}[style=term]
$ dotnet new list

Template Name            Short Name       Language
-----------------------  ---------------  ----------
Console App              console          [C#],F#,VB
Class Library            classlib         [C#],F#,VB
ASP.NET Core Empty       web              [C#],F#
ASP.NET Core Web API     webapi           [C#],F#
xUnit Test Project       xunit            [C#],F#,VB
Solution File            sln
Dotnet gitignore file    gitignore
Dotnet local tool manifest  tool-manifest
EditorConfig file        editorconfig
\end{lstlisting}
  \vspace{0.6em}
  \aside{Не только проекты: \hi{gitignore}, файл решения и манифест
    инструментов создаются отсюда же.}
\end{frame}
\note{Шаблонов десятки. Показать, что часть из них создаёт один файл,
  а не проект.}

\begin{frame}{Жизненный цикл проекта}
  \slidetop
  \claim{От рождения до поставки. Запомним \hi{в формате истории}.}

  \begin{center}
    \begin{tikzpicture}[font=\small,x=2.03cm]
      \draw[hairline,line width=0.6pt] (-0.3,0) -- (6.3,0);
      \foreach \i/\c/\k in {0/new/{-n}, 1/restore/{}, 2/build/{-c},
                            3/run/{--project}, 4/test/{}, 5/add package/{},
                            6/publish/{-c -o}} {
        \fill[muted] (\i,0) circle (1.6pt);
        \node[above=0.45em,font=\small\ttfamily] at (\i,0) {\c};
        \node[below=0.5em,muted,font=\scriptsize\ttfamily] at (\i,0) {\k};
      }
    \end{tikzpicture}
  \end{center}

  \slidegap
  \aside{Проект рождается, обрастает зависимостями, собирается, запускается,
    проверяется и уезжает на сервер. \texttt{restore} --- единственный шаг,
    которому нужна сеть.}
\end{frame}
\note{restore вызывают явно редко: build и run делают его сами.}

\begin{frame}[fragile]{Вся цепочка, три экрана}
  \slidetop
\begin{lstlisting}[style=term]
$ dotnet new console -n Wassup && cd Wassup
$ dotnet run
Hello, World!

$ dotnet run -c Release          # другая конфигурация -> другая папка в bin
Hello, World!

$ dotnet publish -c Release -o ./out
$ ls out
Wassup  Wassup.deps.json  Wassup.dll  Wassup.pdb  Wassup.runtimeconfig.json
\end{lstlisting}
  \vspace{0.6em}
  \aside{В \texttt{out} --- никаких исходников, никакого \texttt{obj}, только
    результат. Это то, что \hi{поедет на сервер}.}
\end{frame}
\note{publish по умолчанию собирает в Release, в отличие от build и run.
  Ключ всё равно пишу явно, чтобы читающий скрипт не гадал.}

\begin{frame}{Инструменты: глобальные и локальные}
  \slidetop
  \versus{глобальные, ключ -g}{%
    Ставятся для вашего пользователя, доступны из любой папки.\\[0.5em]
    {\color{muted}Диагностика: \texttt{dotnet-counters} и подобное.}\\[0.5em]
    Нужен \hi{лично вам}.}{локальные, манифест в репозитории}{%
    Перечислены в манифесте вместе с версиями.\\[0.5em]
    {\color{muted}\texttt{dotnet tool restore} после клонирования.}\\[0.5em]
    Нужен \hi{проекту}.}

  \slidegap
  \aside{Манифест коммитится, в отличие от bin и obj: это исходник,
    а не результат. У глобальных версии разъезжаются между машинами.}
\end{frame}
\note{Я это проходил с инструментом миграций: час на выяснение, что версии
  разные.}

\begin{frame}[plain]
  \keyline{CLI --- не <<сложный способ>>.\\ Это единственный способ, который
    существует везде, где живёт ваше приложение.}
\end{frame}
\note{Сначала команды, потом кнопки. К концу курса кнопки разрешу.}

% ============================================ 5. ПЕРВОЕ ПРИЛОЖЕНИЕ ========
\section{5. Первое приложение}

\begin{frame}[plain]
  \sectionopener{5}{Первое приложение}{Пять файлов, два этапа компиляции,
    ноль магии.}
\end{frame}
\note{Бюджет: 5 минут. Собираем всё в одну картинку.}

\begin{frame}[fragile]{Что было и что стало}
  \slidetop
  \begin{minipage}[t]{0.36\textwidth}
    \eyebrow{после new console}\par\vspace{0.6em}
\begin{lstlisting}[style=plain]
Wassup/
  Wassup.csproj
  Program.cs
\end{lstlisting}
    \vspace{0.5em}
\begin{lstlisting}[style=csharp]
Console.WriteLine("Hello, World!");
\end{lstlisting}
  \end{minipage}\hfill
  \begin{minipage}[t]{0.58\textwidth}
    \eyebrow{после run: bin/Debug/net10.0}\par\vspace{0.6em}
    {\small
    \begin{tabular}{@{}p{3.5cm}p{4.0cm}@{}}
      \texttt{Wassup.dll} & наш код в IL \\[0.4em]
      \texttt{Wassup} & запускатель, нативный \\[0.4em]
      \texttt{Wassup.pdb} & отладочные символы \\[0.4em]
      \texttt{Wassup.deps.json} & от чего зависим \\[0.4em]
      \texttt{Wassup.runtimeconfig.json} & \hi{какой рантайм нужен} \\
    \end{tabular}}
  \end{minipage}

  \vspace{0.9em}
  \aside{Класс и \texttt{Main} тут есть --- их генерирует компилятор.
    \texttt{using} нет, потому что их подставил \texttt{ImplicitUsings}.}
\end{frame}
\note{В лабораторной попросят объяснить именно эти файлы.}

\begin{frame}[fragile]{По этому файлу приложение выбирает рантайм}
  \slidetop
\begin{lstlisting}[style=json]
{
  "runtimeOptions": {
    "tfm": "net10.0",
    "framework": {
      "name": "Microsoft.NETCore.App",
      "version": "10.0.0"
    }
  }
}
\end{lstlisting}
  \vspace{0.7em}
  \aside{Помните \texttt{dotnet -{}-info} в начале лекции --- рантаймов на
    машине было несколько? Вот \hi{по этому файлу} приложение находит свой.}
\end{frame}
\note{Показываю сокращённо: там есть ещё настройки самого рантайма.}

\begin{frame}{Вся цепочка целиком}
  \slidetop
  \begin{center}
    \begin{tikzpicture}[font=\small,
      box/.style={draw=hairline,line width=0.6pt,fill=codebg,
                  minimum height=0.9cm,align=center,inner sep=0.4em},
      ar/.style={-{Straight Barb[length=1.6mm]},muted}]
      \node[box] (src) {Program.cs};
      \node[box,right=0.95cm of src] (dll) {dll с IL\\
        {\scriptsize deps + runtimeconfig + запускатель}};
      \node[box,right=0.95cm of dll] (host) {запускатель\\читает конфиг};
      \node[box,right=0.95cm of host] (rt) {рантайм\\загружает dll};
      \node[box,right=0.95cm of rt] (out) {JIT\\строка на экране};
      \draw[ar] (src) -- (dll);
      \draw[ar] (dll) -- (host);
      \draw[ar] (host) -- (rt);
      \draw[ar] (rt) -- (out);
      \node[muted,font=\scriptsize,below=0.5em] at (src.south) {компиляция};
      \node[muted,font=\scriptsize,below=0.5em] at (rt.south) {запуск};
    \end{tikzpicture}
  \end{center}

  \slidegap
  \aside{Вернёмся к вопросу, с которого начали: что происходит между вашим
    текстом и процессом в операционной системе. Теперь вы можете назвать
    каждый шаг.}
\end{frame}
\note{На первой лабораторной проделаете это сами, целиком из терминала.}

\begin{frame}[plain]
  \statement{Пять файлов, два этапа компиляции, \hi{ноль магии}.}
\end{frame}
\note{Закончить здесь. Не добавлять ничего после.}

\end{document}
